\documentclass[10pt]{article}

% graphicx package, useful for including eps, pdf, and raster graphics
\usepackage{graphicx}
\DeclareGraphicsExtensions{.png,.png,.jpg,.pdf}

% basic packages
\usepackage{fancyhdr}
\usepackage{color}
\usepackage{parskip}
\usepackage{float}
\usepackage{microtype}
\usepackage{url}
\usepackage{hyperref}
\usepackage{booktabs}
\usepackage{makecell}
\usepackage{multirow}
\usepackage{pbox}
\usepackage{enumitem}
\usepackage{xeCJK}

% fonts
\usepackage{fontspec}
\setmainfont{DroidSans}[
  Extension   = .ttf,
  UprightFont = *,
  BoldFont    = *-Bold ]
\setmonofont{DroidSansMono}[
  Extension   = .ttf,
  UprightFont = * ]
\setCJKmainfont{NotoSansCJK}[
  Extension   = .ttc,
  UprightFont = *-Regular,
  BoldFont    = *-Bold,
  Language    = Korean ]

\hypersetup{colorlinks,urlcolor=blue}
\urlstyle{same}

% page layout
\usepackage[top=1in, bottom=1in, left=0.75in, right=0.75in, includefoot, heightrounded]{geometry}

% Useful for seeing the layout as you tweak it.
%\usepackage{showframe}

\setlength{\headheight}{1.1in}
\pagestyle{fancy}
\fancyhf{}

\lhead{\includegraphics[height=1in,clip]{logo}}
\rhead{\pbox[b][1in][c]{\textwidth}{
    \small
    \textbf{검사 수행 실험실:}\\
    Northwest Genomics Center\\
    University of Washington School of Medicine\\
    Genome Sciences Box 355065\\
    3720 15th Ave NE\\
    Seattle, WA 98195-5065}}

\newcommand{\PageLine}{\rule{\textwidth}{0.25mm}}

\begin{document}

\begin{description}[font=\normalfont,align=left,labelwidth=8em]
\item [\textbf{\VAR{pat_name|e}}]
\item [\textbf{생년월일:}] \VAR{birth_date|e}
\item [\textbf{검체:}] 콧속에서 직접 채취
\item [\textbf{채취 검체 바코드:}] \VAR{qrcode|e}
\item [\textbf{채취일:}] \VAR{collect_ts|e}
\item [\textbf{보고일:}]
  %- if status_code == "never-tested"
  테스트한 적이 없음
  %- else
  \VAR{result_ts|e}
  %- endif
\end{description}

\PageLine

\begin{center}
\Large
\textbf{연구 결과 - 의료진과 본 결과를 공유하시기 바랍니다.}
\end{center}

\bigskip

\large \underline{\textbf{검사 관련 정보}}

\begin{itemize}
\item

  COVID-19는 SARS-CoV-2라고 하는 \textbf{바이러스}에 의해 유발되는 \textbf{질환}입니다. 본
  검사에서는 COVID-19 감염 여부를 파악하기 위해 귀하의 코에서 채취한 샘플에 바이러스가
  존재하는지 여부를 파악합니다. 이러한 유형의 검사를 기술적으로 '역전사 중합효소
  연쇄반응(RT-PCR)'이라고 합니다.

\end{itemize}

\bigskip

\large \underline{\textbf{주요 안내사항}}

귀하의 검사 결과를 보기 전에 다음 사항을 숙지해 주세요.

\begin{itemize}
\item

  귀하가 채취한 \textbf{샘플 유형}과 \textbf{실험실 검사}는 \textbf{연구}의 일환이며, 본
  보고서는 귀하의 \textbf{검사 결과}에 대한 설명입니다.

\item

  직접(또는 다른 사람을 대신해서) 샘플을 수집하고 라벨을 붙인 뒤 배송 서비스를 통해 실험실로 회송해
  주세요. 간호사나 의사 등 면허가 있는 전문 의료진이 수집하거나 라벨을 붙이거나 운송하지는
  않습니다. 이는 귀하의 검사 결과에 영향을 미칠 수 있으며, 의료진이 다시 검사를 해야 할 수도
  있습니다.

\item

  본 \textbf{연구 참여}가 의료 서비스를 대체할 수는 없습니다. 많이 아프시거나 검사 결과가
  우려되신다면 의사나 의료진에게 연락하세요.

\end{itemize}

\bigskip
\PageLine

\large \underline{\textbf{SARS-CoV-2(COVID-19) 검사 결과}}

%- if status_code == "negative"
미검출(음성)\\
%- elif status_code == "positive"
검출(양성)\\
%- elif status_code == "inconclusive"
미확정\\
%- elif status_code == "never-tested"
검사 불가\\
%- else
\VAR{0/0} %# poor man's raise
%- endif

\PageLine
\bigskip

\large \underline{\textbf{결과 파악}}

다음 도표는 테스트 결과로 제시될 수 있는 사례와 각각에 대한 설명을 제시한 것입니다. 본
\textbf{검사 결과}를 출력하거나 이메일로 발송함으로써 의사나 의료진과 공유할 수 있습니다.

검사 결과와 관계없이, 아프신 경우에는 자택에 머무르면서 타인과의 접촉을 피하세요. \textbf{증상이
계속되거나 많이 아프거나 질병에 대해 우려가 되거나 검사 결과에 대한 질문이 있으시면 의사나
의료진에게 연락하세요.}

\begin{itemize}


\item

  \textbf{검사 결과 '미검출(음성)'이} 나온 경우에는 실험실에서 귀하가 발송한 샘플을 통해
  COVID-19를 유발하는 바이러스를 찾지 못한 경우에 해당합니다. 그러나, 어떤 검사도 완벽하지는 않기
  때문에 음성 판정을 받았더라도 바이러스가 있을 수 있습니다.

\item

  \textbf{검사 결과 '검출(양성)'이} 나온 경우에는 실험실에서 귀하가 발송한 샘플을 통해
  COVID-19를 유발하는 바이러스를 발견한 경우에 해당합니다. 특히, 아프신 경우에는 자택에 머무르면서
  타인과의 접촉을 피하세요.

  귀하의 결과는 워싱턴 주 및 지역 보건부에도 전달됩니다. 직접 결과를 신고하실 필요는
  \textbf{없습니다}. 보건부가 전화를 통해 귀하의 질환과 귀하가 아픈 동안 접촉했을 수 있는
  사람들에 대해 질문할 수 있습니다.

\item

  \textbf{검사 결과 '미확정'이} 나온 경우에는 실험실이 귀하의 검사 결과를 확신할 수 없는 경우에
  해당됩니다. 어떤 검사도 완벽하지 않기 때문에 샘플을 통해 COVID-19에 감염되었다고 말할 수는
  없지만 소량의 SARS-CoV-2가 발견되는 경우, 미확정에 해당됩니다.

\item

  \textbf{검사 결과 '검사 불가'로} 나온 경우에는 귀하가 제공한 샘플이 부족하여 검사를 할 수
  없거나 검사를 하는 동안 오류가 발생한 경우에 해당됩니다. 이 경우에는 귀하의 샘플을 검사할 수 없기
  때문에 귀하에게 COVID-19를 유발하는 바이러스가 있는지 여부를 파악할 수 없습니다.

\end{itemize}

\bigskip

\large \underline{\textbf{자료}}

COVID-19와 관련된 자세한 정보는 공중 보건국 파트너가 제공하는 다음 자료들을 참조하세요.

\begin{itemize}
\item
  Public Health \textemdash Seattle \& King County(공중 보건국- 시애틀 \& 킹카운티)\\
  \url{https://www.kingcounty.gov/depts/health/covid-19/languages/korean.aspx}

\item

  Washington State Department of Health(워싱턴 주 보건부)\\
  \url{https://www.doh.wa.gov/Emergencies/NovelCoronavirusOutbreak2020COVID19/Korean}

\item
  U.S. Centers for Disease Control and Prevention (CDC, 미 질병통제예방센터)\\
  \url{https://korean.cdc.gov/coronavirus/2019-ncov/index.html}

\end{itemize}

전화 통화를 원하시는 경우에는 다음으로 연락하세요.

\begin{itemize}

\item

  \textbf{COVID-19와 관련된 의학적 질문: } 오전 8시부터 저녁 7시 사이 King County Call
  Center(킹 카운티 콜센터)에 \textbf{206-477-3977} 번으로 연락하세요.

\item

  \textbf{규제 준수 및 비즈니스 관련 문제를 포함하여 COVID-19에 대한 비의학적 질문:} 월요일부터
  금요일 오전 8시 30분부터 오후 4시 30분 사이 King County COVID-19 Business and
  Community Information Line(킹 카운티 COVID-19 비즈니스 및 지역사회 정보 라인)에
  \textbf{206-296-1608}

\item

  \textbf{워싱턴 주 COVID-19에 대한 일반 문의:} 오전 6시부터 밤 10시 사이 Washington
  State Novel Coronavirus Call Center(워싱턴 주 신종 코로나바이러스 콜센터)에
  \textbf{1-800-525-0127} 번으로 연락하세요.

\end{itemize}

\bigskip

\large \underline{\textbf{방법 및 제한사항}}

\textbf{방법}\\
본 검정은 역전사 PCR을 사용하여 SARs-CoV-2를 검출하기 위한 것입니다. 본 검정은 바이러스 세그먼트
ORf1B 및 S의 증폭을 이용하며 인간 RnaseP 증폭 통제도 검출합니다. University of
Washington의 Northwest Genomics Center가 본 검사를 개발하고 그 성능 특성을 파악했습니다.
본 실험실은 고도로 복잡한 임상실험검사를 수행할 수 있는 자격을 갖춘 WA Department of Health
의학검사현장으로 허가받았습니다.

\textbf{제한사항}

\begin{itemize}

\item

  본 결과는 연구의 일환으로 제공됩니다.

\item

  직접(또는 다른 사람을 대신해서) 샘플을 수집하고 라벨을 붙인 뒤 배송 서비스를 통해 실험실로 회송해
  주세요. 간호사나 의사 등의 전문 의료진이 수집하거나 라벨을 붙이거나 운송하지는 않습니다. 이는
  귀하의 검사 결과에 영향을 미칠 수 있으며, 의료진이 다시 검사를 해야 할 수도 있습니다.

\item

  본 검사는 연구 데이터를 토대로 하기 때문에, 귀하의 결과로 SARS-CoV-2 감염에 의해 유발되는
  질환인 COVID-19 감염 여부를 진단하지 못할 수 있습니다.

\end{itemize}

\end{document}
